\section*{Part III: Price Metrics and Commercial Architecture}

\subsection*{15. What Unit Is the Buyer Actually Purchasing?}

A price requires a denominator: per unit, per kilogram, per seat, per learner, per cohort, per project, per month, per transaction, per call, per successful job, per license, or per percentage of attributable value. The unit determines how the customer's bill changes with use, and often how the customer behaves.

An effective pricing metric is intelligible, measurable, resistant to gaming, aligned with buyer value, and compatible with seller cost. It should not punish desired adoption unnecessarily. A usage metric may be fair when value and variable cost both grow with usage; it can be destructive when customers avoid using the tool to control an unpredictable bill. A per-seat metric may align with enterprise budgets while discouraging collaboration.

Pricing architecture is a hypothesis about how value scales. If the unit does not track value, the customer and seller experience friction as adoption grows.

\subsection*{16. Fixed Prices and Per-Unit Sales}

A fixed unit price reduces explanation burden and helps customers compare alternatives. It works naturally for many consumer goods, standardized diagnostic services, and clearly scoped workshops. The risk is customer heterogeneity: heavy users may cost more; others may value only a subset of the offer. 1 price may exclude people whose job would be better served by a simpler experience.

For physical goods, the founder must include packaging, logistics, retailer economics, inventory loss, and returns; the price on a shelf differs from the manufacturer's realized revenue. For services, the exact scope and change conditions matter. ``Fixed fee'' without ``fixed specification'' often becomes an unlimited work obligation.

\subsection*{17. Hourly Fees, Project Fees, and Retainers}

Hourly billing connects fees to labor and makes scope changes visible. It is sometimes appropriate for uncertain legal, technical, or advisory work. It can also penalize expert efficiency and encourage disputes over how time is used. A project fee transfers execution risk to the provider and works better when deliverables, dependencies, acceptance criteria, change orders, and exclusions are explicit.

Retainers sell ongoing availability or a defined service envelope. They help the customer access expertise predictably and the seller plan capacity. They fail when ``access'' implies unlimited work, when buyers feel they must use time to justify the fee, or when the provider oversells availability. The real product is not hours; it is a clearly governed professional commitment.

\subsection*{18. Subscription: Recurring Billing Needs Recurring Value}

Subscriptions can suit software, maintenance, information services, memberships, and repeated assistance. They smooth bills for customers and support planning for sellers. But recurring billing is not itself recurring customer value. A 1-time learning transformation should not be forced into a monthly subscription simply because subscriptions look attractive in investor presentations.

A subscription requires a clear ongoing job: updates, access, collaboration, compliance, monitoring, coaching, or regular content that matters. The model must track active use, renewal decision, churn, service obligations, support cost, and the point at which accumulated value becomes greater than the cost of switching. A customer who continues paying from inertia may inflate short-run revenue but damage trust once they notice.

Annual prepayment improves cash at the cost of future delivery obligations. Revenue is generally recognized as services are transferred under applicable rules, not because the cash arrived. A founder should know exactly what has been earned and what remains owed.

\subsection*{19. Usage Pricing and AI's Variable Marginal Costs}

AI inference, storage, transactions, and compute can have meaningful variable cost. A product that bills a flat fee to all users may subsidize exceptionally heavy workloads. Usage charges can better align cost and customer benefit, but output tokens, API requests, and compute minutes are often poor proxies for the \emph{success} customers want.

A model-call pricing metric can reward customers for doing less work with the product, while a per-completed-case metric encourages reliability. The latter is harder to define: what counts as successfully completed, how are errors resolved, and who verifies outcomes? Hybrid pricing: base fee plus usage: can cover availability while allowing expansion. Model potential bill shock, spend caps, and quality loss before launching.
